\chapter*{Реферат}
\addcontentsline{toc}{chapter}{Реферат} 

\begin{center}
    Общая характеристика диссертации
\end{center}

%\begin{figure}
%    \centering
%    \includegraphics[width=0.6\linewidth]{images/knuth}
%    \caption{Knuth}
%    \label{fig:my_label2}
%\end{figure}


\paragraph*{Актуальность темы исследования.}

\paragraph*{Степень разработанности темы.}

\paragraph*{Цель исследования.}
увеличение эффективности размещения объектов обслуживания населения в рамках процесса территориального планирования

\paragraph*{Задачи исследования.}
Исследование системы управления территориальным планированием в задаче размещения объектов социального обслуживания населения;
Разработка метода оценки транспортной доступности на основе критерия сезонной устойчивости;
Разработка метода оптимизации сети в задаче оптимального размещения;
Применение методов на примере агломераций Арктической зоны РФ.

\paragraph*{Объект исследования.}
подсистема территориального планирования в составе системы управления развитием территорий

\paragraph*{Предмет исследования.}
методы и алгоритмы поддержки принятия решений по обеспечению транспортной доступности объектов обслуживания

\paragraph*{Методы исследования.}

\paragraph*{Научная новизна.}

\paragraph*{Теоретическая значимость.}

\paragraph*{Практическая значимость.}

\paragraph*{Положения, выносимые на защиту.}
\begin{enumerate}
    \item Метод оценки транспортной доступности на основе критерия сезонной устойчивости, отличающийся применением методов сетевого анализа, обеспечивающий учет стохастической динамики многоуровневой сети. Положение соответствует пункту 2 паспорта специальности.
    
    \item Метод интеллектуальной поддержки принятия решений оптимального размещения объектов обслуживания, отличающийся применением оптимизации сети, обеспечивающий уменьшение необходимого количества объектов обслуживания. Положение соответствует пункту 9 паспорта специальности.
\end{enumerate}

\paragraph*{Достоверность результатов.}
Степень достоверности научных достижений
подтверждается корректным использованием методов, обоснованием постановки задач, экспериментальными исследованиями на реальных массивах данных, покрывающими разработанные технологии и алгоритмы.

\paragraph*{Апробация результатов.}
Результаты исследований обсуждались на конференциях:​

\begin{enumerate}
	\item 40-я ежегодная конференция AAAI по искусственному интеллекту (2026) --- 2-й семинар по ИИ для городского планирования, 2026, рейтинг CORE: A*, стендовый и устный доклад;
	\item Научная конференция AMS (2026) --- признанная площадка исследований в области городского планирования, специальная секция, устный доклад;
	\item Международная школа и конференция по науке о сетях (2026) --- ведущая площадка исследований в области сетевой науки, основная программа, стендовый доклад;
	\item Международная конференция по вычислительной науке и ее приложениям (ICCSA 2023/2024).
\end{enumerate}

\paragraph*{Публикации автора по теме диссертации.}
Основные результаты по теме диссертации изложены в 3 публикациях,
опубликованных в изданиях, индексируемых в базах цитирования Scopus и Web of Science.


\paragraph*{Личный вклад автора.}
В работах, выполненных в соавторстве, вклад автора заключается в
построении моделей, разработке и реализации методов и алгоритмов, разработке, проведении вычислительных экспериментов и подготовке обзоров литературы. 

В совместной статье с Губиной А. автор выполнил обзор методов моделирования многоуровневых сетей, обеспечил программную реализацию и разработал алгоритмы расчетов доступности при сезонном изменении в сети.

Во всех работах, выполненных в соавторстве с Тихевич В., автор выполнил постановку задач, методологический обзор и разработку предложенных алгоритмов.

В работах, выполненных совместно с Морозовым А.С. автор выполнил методологическую постановку экспериментов, программную реализацию, а также анализ полученных результатов.



\paragraph*{Объём и структура работы.}
Диссертация состоит из~введения,
\formbytotal{totalchapter}{глав}{ы}{}{},
заключения и
\formbytotal{totalappendix}{приложен}{ия}{ий}{}.
%% на случай ошибок оставляю исходный кусок на месте, закомментированным
%Полный объём диссертации составляет  \ref*{TotPages}~страницу
%с~\totalfigures{}~рисунками и~\totaltables{}~таблицами. Список литературы
%содержит \total{citenum}~наименований.
%
Полный объём диссертации составляет
\formbytotal{TotPages}{страниц}{у}{ы}{}, включая
\formbytotal{totalcount@figure}{рисун}{ок}{ка}{ков} и
\formbytotal{totalcount@table}{таблиц}{у}{ы}{}.
Список литературы содержит
\formbytotal{citenum}{наименован}{ие}{ия}{ий}.

\paragraph*{Соответствие специальности 2.3.4 (в порядке степени соответствия).}

\begin{itemize}
	\item[2.] Разработка математических моделей и критериев эффективности, качества и надёжности организационных систем;
	\item[9.] Разработка методов и алгоритмов интеллектуальной поддержки принятия управленческих решений в организационных системах.
\end{itemize}


\newpage
\section*{Основное содержание работы}

В Главе~\ref{ch:ch1}...


